\section{本章小结}
\label{sec:ls-equations-summary}

本章研究已知输出、反求输入的问题：先是线性方程$\bm A\bm x=\bm b$，再是以矩阵为未知量的Sylvester方程$\bm A\bm X+\bm X\bm B=\bm C$与Lyapunov方程。线性方程部分依次回答了四个问题：要求什么样的解，怎样算出来，结果能相信多少，以及怎样选择方法。

求解目标由秩决定。$\operatorname{rank}[\bm A\;\bm b]=\operatorname{rank}(\bm A)$时方程有解，全部解是一个特解加上$\operatorname{null}(\bm A)$；无解时改求最小二乘解，其残差与列空间正交，满足总是相容的正规方程；解不唯一时选落在$\operatorname{range}(\bm A^\top)$中的最小范数解。伪逆$\bm A^\dagger\bm b$把三种情况的选解写成同一个公式。加权最小二乘经白化化为普通情形；齐次最小二乘在单位长度下最小化$\lVert\bm A\bm x\rVert_2$，解是最小右奇异向量。

计算路线有四类。直接法先分解再回代：一般方阵用带部分选主元的LU，对称正定用Cholesky，对称不定用分块LDL分解，分块结构用Schur补，低秩修改用Woodbury公式复用已有分解；最小二乘用正规方程或QR，其中Householder与Givens以正交变换保持长度，修正Gram--Schmidt在矩阵病态时可能失去正交性。定常迭代把矩阵分裂为易解部分与余项，Richardson、Jacobi、Gauss--Seidel、SOR与SSOR的收敛都由迭代矩阵的谱半径决定。Krylov子空间法在$\bm x_0+\mathcal K_k$中按不同准则选取近似：CG使误差的$\bm A$范数最小，GMRES与MINRES使残差最小，BiCG与BiCGSTAB用短递推处理非对称方程；LSQR与LSMR基于Golub--Kahan双对角化，在数学上分别等价于对正规方程施行CG与MINRES，但不形成$\bm A^\top\bm A$；预条件改变的是收敛速度，有时也改变所求的解。Kaczmarz每步只用一行，相容时收敛，扩展版本再用列投影去掉右端不可拟合的部分，从而收敛到最小二乘解。

结果的可信程度由数据误差和计算误差共同决定。条件数刻画解对数据扰动的放大，残差小不等于误差小；后向误差把计算结果解释为相近问题的精确解，前向误差约不超过条件数乘以后向误差。形成正规方程会把条件数平方；伪逆在秩变化处不连续，所以数值秩要按误差水平判断。岭回归与截断SVD改变问题本身，用拟合误差换取对小奇异值方向的稳定。选择方法时，先固定求解目标，再按矩阵结构、数据访问方式、计算与存储开销、收敛速度和稳定性比较。

Sylvester方程对任意右端唯一可解，当且仅当$\bm A$与$-\bm B$没有公共特征值。求解时不形成Kronecker系数矩阵，而是对$\bm A$、$\bm B$作实Schur分解，逐列解移位的拟三角方程，再变换回原坐标；可对角化时换到特征坐标虽然简单，但特征向量病态时会放大误差，所以一般仍用Schur路线。Lyapunov方程是对称特例：右端对称时解也对称，求解只需算一半未知量。离散型方程在$\rho(\bm A)<1$时有级数解，由此读出右端正定时解正定。
